% =====================================================================
% PARTE II — Revisão crítica do estado da arte  [~10 min]
% Fonte: Chapters/tsad, Chapters/reconstruction, Chapters/forecasting
% =====================================================================

\section{Estado da Arte}

\subsection{Detecção de anomalias}

\begin{frame}{Taxonomia de anomalias em séries temporais}
  \duascolunas{%
    \begin{itemize}
      \item Pontuais
      \item Contextuais
      \item Coletivas
    \end{itemize}
    \vspace{0.5em}
    {\scriptsize\color{peegray}Critérios em prosa; sem equações ilustrativas.}
  }{%
    \centering
    \pendente{usar \texttt{anomaly\_taxonomy.png} do workspace}
  }
\end{frame}

\begin{frame}{Famílias de métodos e o problema da avaliação}
  \begin{itemize}
    \item Distância, densidade e predição; realizações neurais.
    \item \dest{Benchmarks} e o viés de \emph{point-adjust}.
    \item Medidas de localização: Tatbul, VUS, afiliação, PATE.
    \item \pendente{fechar com a lacuna que justifica a proposta}
  \end{itemize}
\end{frame}

\subsection{Reconstrução e previsão}

\begin{frame}{Reconstrução: contrato da tarefa}
  \begin{itemize}
    \item Três estatutos do alvo.
      \pendente{Cap. 5}
    \item Reconstrução \alerta{$\neq$} correção de anomalia.
  \end{itemize}
\end{frame}

\begin{frame}{Previsão: do ARIMA aos Transformers}
  \begin{itemize}
    \item Métodos clássicos, aprendizado de máquina e profundo.
    \item Transformers para séries temporais: ganho real e custo.
    \item \pendente{Cap. 6}
  \end{itemize}
\end{frame}

\begin{frame}{Síntese: as três lacunas}
  \begin{enumerate}
    \item \pendente{lacuna 1}
    \item \pendente{lacuna 2}
    \item \pendente{lacuna 3}
  \end{enumerate}
  \vfill
  \begin{block}{Consequência}
    \pendente{a transição para a proposta}
  \end{block}
\end{frame}
